% \begin{frame}[t]{Freiburg Ion Trapping Group}
%     \picpos{\linewidth}{0cm}{0cm}{group_slide.pdf}
% \end{frame}

\begin{frame}[t]{BaLi: An Introduction}
  \visible<1->{
  Why \BaLi?
  \begin{itemize}
    \item Looking for interactions in the ultracold regime
    \item<1-> Large mass ratio means lower collision energy
    \item<1-> Chemically stable in their ground states
    \item<1-> Li can be either bosonic or fermionic
    \item<1-> Candidates for all-optical trapping
    \end{itemize}
  }

\end{frame}

\begin{frame}[t]{Atom-Ion interaction}
    \only<1>{The ion induces a dipole moment in the atom ($p \propto E(r) \propto \frac{1}{r^2}$) that leads to a long-range charge-dipole interaction:}

    \only<2>{The total wavefunction is composed of partial waves of angular momentum $l$.

    The contribution of partial waves depends on the collision energy.}
    \[
        \setlength{\abovedisplayskip}{2pt}
         \setlength{\belowdisplayskip}{2pt}
        V_\text{atom-ion}(r) \approx- \frac{C_4}{r^4} \visible<2>{\textcolor{red}{+ \frac{\hbar \,l\,(l+1)}{\mu r^2}}}
        \]
        \only<1>{\picpos{0.4\linewidth}{.5\linewidth}{0cm}{BaLi_PECs_B.pdf}
        \picpos{0.35\linewidth}{1cm}{0cm}{atom_ion_potential.png}}
        \only<2>{\picpos{.75\linewidth}{0.1\linewidth}{0cm}{partial_waves.png}}
\end{frame}


\begin{frame}[t]{The Feshbach Resonances}
    \only<1-2>{An incoming state (open channel) can couple to a bound state (closed channel).

    As the bound state approaches the incoming state the mixing between states increases.

    \picpos{.55\linewidth}{0cm}{0.6cm}{feshbach_resonance2.png}}
  %   \only<1>{\begin{textblock*}{0.4\linewidth}(0.5\linewidth,0cm)
  %   Quantum number $M=m_1 + m_2 + m_l$ is strictly conserved.
  % \end{textblock*}}
    \only<2>{
        \picpos{.46\linewidth}{.53\linewidth}{2.4cm}{feshbach_resonance_isolated.pdf}
        \begin{textblock*}{.5\linewidth}(.51\linewidth, 0.2cm)
           \small
           Magnetic fields can be used to shift potentials relative to each other (magnetic FR)
                   \[ a = a_\text{bg} \left( 1 - \frac{\Delta}{B-B_0} \right) \]
        \end{textblock*}
   }
\end{frame}

\begin{frame}[t]{Feshbach Resonances}
    Why are we interested in Feshbach resonances?
    \begin{itemize}
      \item Molecular spectroscopy
      \item Tuning of interaction rates
      \item Interaction can be made attractive ($a>0$), repulsive ($a<0$) or switched off ($a=0$)
      \item Formation of cold molecules in chosen states
      \end{itemize}
\end{frame}
